\begin{frame}{리뷰 템플릿 (5번의 두 칸이 핵심)}
  \scriptsize
  \begin{tabular}{@{}p{0.16\textwidth}p{0.8\textwidth}@{}}
    \toprule
    칸 & 내용 \\
    \midrule
    1. 환경 설명 & 상태·행동·보상·주 지표(greedy 평가 리턴? 학습중 리턴?)가
      보고서에서 추적 가능한가? \textbf{베이스라인(무작위) 리턴}은
      명시됐는가? \\
    2. 알고리즘 선택 & 이 환경의 어떤 특성(상태 수, 보상 구조, 확률 전이)
      이 비교를 정당화하는가? 그 연결이 Ch\,4$\sim$7 개념으로
      설명되는가? \\
    3. 수식적 정당화 & ``$\geq 1$개''(8.1)를 채웠는가? 짚은 개념이
      \textbf{실제로 관측한 숫자}와 일치하는가(단순 인용이 아닌가)? \\
    4. 결과의 정직성 & 시드 $\geq 3$, 시드별 숫자 표, 평가 프로토콜
      ($\varepsilon=0$? 에피소드 수?) 명시. 불안정한 시드·구간이
      \textbf{정직하게} 보고됐는가? \\
    5. \textbf{반박 가능한 질문(필수, 1개 이상)} &
      \textbf{근거가 되는 개념}: Week \_\_ 절 \_\_ (개념: \_\_\_) \quad
      \textbf{발표팀이 답하려면 필요한 새 실험/계산}: \_\_\_ \quad
      \textbf{질문}: \_\_\_ (발표팀이 \emph{답이 나쁘면} 결과가
      흔들리는 것이어야) \\
    \bottomrule
  \end{tabular}
  \vskip0.6em
  \begin{block}{3점 도달 조건}
    5번의 ``개념''과 ``새 실험'' \textbf{두 칸이 모두} 채워져야 3점.
    둘 중 하나가 비어 있으면 \textbf{반박 불가능}한 리뷰.
  \end{block}
\end{frame}
